\scanpage{78}
\mspage{84}

Quand ces conditions sont satisfaites, on s'attend à ce que le foncteur
\[
f^*:\mathbb D(Y)\longrightarrow\mathbb D(X)
\]
soit pl. fid, i.e. que les morphismes d'adjonction
\[
F\unsure{-|\longrightarrow} \unsure{p_*|f_*|f} f^*F,\qquad f_! f^*\unsure{f|F}\unsure{-|\longrightarrow} F
\]
dans $\mathbb D(Y)$ soient des iso. (C'est vrai si $f$ est $\mathbb D$-propre ou $\mathbb D$-lisse, cf cor. 3 p.~71.) Cela revient encore à dire que les $p_*\unsure{f^*|p^*}F$ se calculent par passage aux fibres, \deleted{\unsure{et}} \inserted{ou encore} itou pour les $p_!p^*F$. Cela nous ramène à la question de la page \unsure{87|82|89},

\uline{Définition 3} On dit alors que $f$ est une $\mathbb D$ fibration \uline{triviale}.

\uline{Proposition} 6 (précise prop. 3 page 81) Considérons un triangle comm. dans $\mathrm{Diag}$
\diagram{\begin{tikzcd}[ampersand replacement=\&]
X\arrow[rr,"f"]\arrow[dr,"p"'] \& \& X'\arrow[dl,"p'"]\\
\& S\&
\end{tikzcd}}

a) \deleted{\unsure{b)}} Supposons que $p$ et $p'$ soient des $\mathbb D$-fibrations gauches, ou que $p$ et $p'$ soient des $\mathbb D$-fibrations droites. Soit \deleted{\unsure{$T$}} $S_0$ une partie de $S$ qui rencontre toute composante connexe de $S$. Pour que $f$ soit une $\mathbb D$-équivalence, il \uline{suffit} que $\forall s_0\in S_0$, $f$ induise une $\mathbb D$-équivalence $X_{s_0}\unsure{-|\longrightarrow} Y_{s_0}$.

% slanted
\margin{cf contre-exemple page 87}

b) Cette condition est aussi nécessaire \inserted{(?)} si $p,p'$ sont des $\mathbb D$-fibrations. Donc dans ce cas, $f$ est une $\mathbb D$-équivalence ssi $f$ induit des $\mathbb D$-équivalences fibre par fibre.

% slanted
\margin{? \uuline{Faux} Pas prouvé ! Même si $X'=S$ ce n'est pas clair, même $\mathbb D=\mathrm{HOT}_{\unsure{\mathrm{AB}|\mathrm{Ab}}}$, $f$ propre et lisse.}

% slanted
\margin{Mais OK. si $S$ $\unsure{\mathbb D\text{-contractile}|\text{contractile}}$\ldots}
% archive sheet number: 77
